\documentclass[11pt,a4paper]{article}
\usepackage[utf8]{inputenc}
\usepackage[T1]{fontenc}
\usepackage[spanish,es-minimal,es-noshorthands]{babel}
\AtBeginDocument{\renewcommand{\sin}{\operatorname{sen}}\renewcommand{\arcsin}{\operatorname{arcsen}}}
\usepackage{lmodern}
\usepackage[sfdefault,lining]{FiraSans}
\usepackage[a4paper,top=2.3cm,bottom=2.2cm,left=1.8cm,right=1.8cm,headheight=1.3cm,headsep=0.4cm,footskip=1cm]{geometry}
\usepackage{amsmath,amssymb}
\usepackage{icomma}
\usepackage{graphicx}
\usepackage[table]{xcolor}
\usepackage{tikz}
\usepackage{fontawesome5}
\usepackage{ccicons}
\usepackage{fancyhdr}
\usepackage{lastpage}
\usepackage{tasks}
\usepackage{needspace}
\usepackage[most]{tcolorbox}
\usepackage{array}
\usepackage[hidelinks]{hyperref}
\AtBeginDocument{\hypersetup{pdftitle={Ficha de ejercicios de la unidad 9: Combinatoria y probabilidad. 1º Bach.},pdfsubject={Matemáticas CCSS 1º Bach.},pdfkeywords={pdfkeywords={combinatoria, recuento, variaciones, permutaciones, combinaciones, números combinatorios, probabilidad, regla de Laplace, probabilidad condicionada, sucesos independientes, experiencias compuestas, probabilidad total, teorema de Bayes}}}}

\definecolor{navy}{HTML}{1B2748}
\definecolor{accent}{HTML}{E8603A}
\definecolor{teal}{HTML}{1F9E90}
\definecolor{soft}{HTML}{EEF1F8}
\definecolor{muted}{HTML}{5D6685}

\setlength{\parindent}{0pt}
\setlength{\parskip}{3pt}
\settasks{label=\alph*), label-format=\bfseries\color{navy}, label-width=1.3em, item-indent=1.8em, column-sep=0.6em, after-item-skip=5pt, before-skip=3pt, after-skip=4pt}

% ---------- capçalera i peu ----------
\newcommand{\llicencia}{\href{https://creativecommons.org/licenses/by-nc-sa/4.0/deed.es}{\ccbyncsa}\ \ Licencia Creative Commons Atribución-NoComercial-CompartirIgual 4.0 Internacional (CC BY-NC-SA 4.0)}
\pagestyle{fancy}
\fancyhf{}
\fancyhead[L]{\raisebox{-0.35\height}{\includegraphics[height=0.95cm]{logo-gva.png}}}
\fancyhead[C]{\small\color{muted}Ficha de ejercicios · Unidad 9: Combinatoria y probabilidad · 1\textsuperscript{o} Bach.}
\fancyhead[R]{\raisebox{-0.35\height}{\includegraphics[height=0.95cm]{logo-vermellar.png}}}
\fancyfoot[L]{\scriptsize\color{muted}\llicencia}
\fancyfoot[R]{\small\color{muted}\thepage}
\renewcommand{\headrulewidth}{0.4pt}
\renewcommand{\headrule}{\color{muted!40}\hrule height \headrulewidth}
\fancypagestyle{primera}{\fancyhf{}\renewcommand{\headrulewidth}{0pt}
  \fancyfoot[L]{\scriptsize\color{muted}\llicencia}
  \fancyfoot[R]{\small\color{muted}\thepage}}

% ---------- nivells i exercicis ----------
\newcommand{\estrelles}[1]{{\color{accent}\ifcase#1\or\faStar\faStar[regular]\faStar[regular]\or\faStar\faStar\faStar[regular]\or\faStar\faStar\faStar\fi}}
\newcounter{exer}
\newcommand{\ex}[1]{\par\needspace{4\baselineskip}\medskip\refstepcounter{exer}\noindent{\color{navy}\bfseries\theexer.}\ \ }

\newcommand{\seccio}[2]{\par\needspace{16\baselineskip}\bigskip
  \begin{tcolorbox}[colback=navy,colframe=navy,coltext=white,arc=2mm,boxrule=0pt,left=3mm,right=3mm,top=1.2mm,bottom=1.2mm]
  {\large\bfseries #1\quad #2}
  \end{tcolorbox}\vspace{-2pt}}

\usetikzlibrary{calc,trees}
\usepackage{enumitem}
\newcommand{\exf}[1]{\par\medskip\refstepcounter{exer}\noindent{\color{navy}\bfseries\theexer.}\ \ }
\providecommand{\gt}{>}\providecommand{\lt}{<}
% quadre de recordatori (només a la fitxa d'exercicis)
\newcommand{\recorda}[1]{\par\smallskip\begin{tcolorbox}[colback=accent!7,colframe=accent!55,arc=2mm,boxrule=0.5pt,left=2.5mm,right=2.5mm,top=1mm,bottom=1mm]\small{\color{accent}\faLightbulb[regular]}\ \textbf{Recuerda:}\ #1\end{tcolorbox}}
\newcommand{\res}[1]{#1}\newcommand{\rest}[1]{#1}

\def\fmneed{11}% que el títol d'una secció no quede sol al final del full
% =====================================================================
% Capçalera FentMates de les fitxes: un color per a cada curs.
% Els colors són els de la pàgina inicial (LEVELS a assets/js/materials.js);
% el logo ve de _eines/logo/pdf/ (python3 _eines/logo/genera.py).
%
% Ús, al final del preàmbul de la unitat:
%   % =====================================================================
% Capçalera FentMates de les fitxes: un color per a cada curs.
% Els colors són els de la pàgina inicial (LEVELS a assets/js/materials.js);
% el logo ve de _eines/logo/pdf/ (python3 _eines/logo/genera.py).
%
% Ús, al final del preàmbul de la unitat:
%   % =====================================================================
% Capçalera FentMates de les fitxes: un color per a cada curs.
% Els colors són els de la pàgina inicial (LEVELS a assets/js/materials.js);
% el logo ve de _eines/logo/pdf/ (python3 _eines/logo/genera.py).
%
% Ús, al final del preàmbul de la unitat:
%   \input{../comu/fentmates.tex}
%   \fmcurs{1eso}
% i, al començament del document:
%   \capcalera{Matemàtiques}{1r ESO}{Fitxa d'exercicis · Unitat 2}{Divisibilitat}
% Opcional abans de \input: \def\fmneed{14} (espai mínim abans de cada secció).
% =====================================================================
\makeatletter
\def\fm@logodir{../../logo/pdf}
\providecommand{\fmneed}{10}
\def\fm@cursnom{}

% clau del curs, color i color del text damunt del color
\newcommand{\fm@curs}[3]{\@namedef{fm@color@#1}{#2}\@namedef{fm@text@#1}{#3}}
\fm@curs{1eso}{E6B84A}{navy}
\fm@curs{2eso}{3B82F6}{white}
\fm@curs{3eso}{FF7A50}{white}
\fm@curs{4eso-a}{A78BFA}{white}
\fm@curs{4eso-b}{6D28D9}{white}
\fm@curs{1batx-mi}{9B1C1C}{white}
\fm@curs{1batx-ccss}{EF6B6B}{white}
\fm@curs{2batx-mii}{1E6B47}{white}
\fm@curs{2batx-ccss}{2BB3A3}{white}

\newcommand{\fmcurs}[1]{%
  \@ifundefined{fm@color@#1}{\PackageError{fentmates}{Curs desconegut: #1}{}}{}%
  \edef\fm@tmp{\noexpand\definecolor{curs}{HTML}{\@nameuse{fm@color@#1}}%
               \noexpand\colorlet{curstext}{\@nameuse{fm@text@#1}}}\fm@tmp
  \colorlet{cursfosc}{curs!72!black}%
  \def\fm@logo{\fm@logodir/vertical-#1.pdf}%
  \def\fm@marca{\fm@logodir/marca-#1.pdf}}

% ---------- capçalera de la primera pàgina ----------
% #1 matèria (petit), #2 curs, #3 tipus de document i unitat, #4 títol
\newcommand{\capcalera}[4]{%
  \gdef\fm@cursnom{#2}%
  \par\noindent
  \begin{tcolorbox}[colback=curs,colframe=curs,boxrule=0pt,arc=2.2mm,width=3.5cm,height=2.15cm,
      valign=center,left=3mm,right=2mm,top=1mm,bottom=1mm,nobeforeafter,box align=center]
    \color{curstext}{\footnotesize #1}\par\vspace{1pt}{\fontsize{19}{21}\selectfont\bfseries #2\par}%
  \end{tcolorbox}\hspace{5mm}%
  \begin{minipage}[c]{\dimexpr\linewidth-3.5cm-5mm-2.9cm\relax}
    {\color{cursfosc}\small\bfseries\MakeUppercase{#3}}\par\vspace{2pt}
    {\color{navy}\fontsize{20}{23}\selectfont\bfseries #4\par}
  \end{minipage}\hfill
  \begin{minipage}[c]{2.6cm}\raggedleft\includegraphics[width=2.6cm]{\fm@logo}\end{minipage}\par
  \vspace{2pt}}

% ---------- barres de secció del color del curs ----------
\renewcommand{\seccio}[2]{\par\needspace{\fmneed\baselineskip}\bigskip
  \begin{tcolorbox}[colback=curs,colframe=curs,coltext=curstext,arc=2mm,boxrule=0pt,left=3mm,right=3mm,top=1.2mm,bottom=1.2mm]
  {\large\bfseries #1\quad #2}
  \end{tcolorbox}\vspace{-2pt}}

% ---------- estrelles de dificultat del color del curs (a Batxillerat no se'n posen) ----------
\@ifundefined{estrelles}{}{%
  \renewcommand{\estrelles}[1]{{\color{curs}\ifcase#1\or\faStar\faStar[regular]\faStar[regular]\or\faStar\faStar\faStar[regular]\or\faStar\faStar\faStar\fi}}}

% ---------- capçalera i peu de les pàgines ----------
% primera pàgina: logos institucionals al peu
\fancypagestyle{primera}{\fancyhf{}\renewcommand{\headrulewidth}{0pt}
  \fancyfoot[L]{\raisebox{-0.32\height}{\includegraphics[height=0.78cm]{logo-gva.png}}\hspace{5mm}%
    \raisebox{-0.32\height}{\includegraphics[height=0.78cm]{logo-vermellar.png}}}
  \fancyfoot[R]{\scriptsize\color{muted}\href{https://creativecommons.org/licenses/by-nc-sa/4.0/deed.ca}{\ccbyncsa}\ \ CC BY-NC-SA 4.0 · fentmates.org\qquad\small\thepage}}
% la resta: el símbol a l'esquerra i el curs a la dreta, amb el títol al centre com abans
\fancyhead[L]{\raisebox{-0.3\height}{\includegraphics[height=0.62cm]{\fm@marca}}}
\fancyhead[R]{\small\bfseries\color{cursfosc}\fm@cursnom}
\renewcommand{\headrule}{\color{curs!45}\hrule height \headrulewidth}
\makeatother

%   \fmcurs{1eso}
% i, al començament del document:
%   \capcalera{Matemàtiques}{1r ESO}{Fitxa d'exercicis · Unitat 2}{Divisibilitat}
% Opcional abans de \input: \def\fmneed{14} (espai mínim abans de cada secció).
% =====================================================================
\makeatletter
\def\fm@logodir{../../logo/pdf}
\providecommand{\fmneed}{10}
\def\fm@cursnom{}

% clau del curs, color i color del text damunt del color
\newcommand{\fm@curs}[3]{\@namedef{fm@color@#1}{#2}\@namedef{fm@text@#1}{#3}}
\fm@curs{1eso}{E6B84A}{navy}
\fm@curs{2eso}{3B82F6}{white}
\fm@curs{3eso}{FF7A50}{white}
\fm@curs{4eso-a}{A78BFA}{white}
\fm@curs{4eso-b}{6D28D9}{white}
\fm@curs{1batx-mi}{9B1C1C}{white}
\fm@curs{1batx-ccss}{EF6B6B}{white}
\fm@curs{2batx-mii}{1E6B47}{white}
\fm@curs{2batx-ccss}{2BB3A3}{white}

\newcommand{\fmcurs}[1]{%
  \@ifundefined{fm@color@#1}{\PackageError{fentmates}{Curs desconegut: #1}{}}{}%
  \edef\fm@tmp{\noexpand\definecolor{curs}{HTML}{\@nameuse{fm@color@#1}}%
               \noexpand\colorlet{curstext}{\@nameuse{fm@text@#1}}}\fm@tmp
  \colorlet{cursfosc}{curs!72!black}%
  \def\fm@logo{\fm@logodir/vertical-#1.pdf}%
  \def\fm@marca{\fm@logodir/marca-#1.pdf}}

% ---------- capçalera de la primera pàgina ----------
% #1 matèria (petit), #2 curs, #3 tipus de document i unitat, #4 títol
\newcommand{\capcalera}[4]{%
  \gdef\fm@cursnom{#2}%
  \par\noindent
  \begin{tcolorbox}[colback=curs,colframe=curs,boxrule=0pt,arc=2.2mm,width=3.5cm,height=2.15cm,
      valign=center,left=3mm,right=2mm,top=1mm,bottom=1mm,nobeforeafter,box align=center]
    \color{curstext}{\footnotesize #1}\par\vspace{1pt}{\fontsize{19}{21}\selectfont\bfseries #2\par}%
  \end{tcolorbox}\hspace{5mm}%
  \begin{minipage}[c]{\dimexpr\linewidth-3.5cm-5mm-2.9cm\relax}
    {\color{cursfosc}\small\bfseries\MakeUppercase{#3}}\par\vspace{2pt}
    {\color{navy}\fontsize{20}{23}\selectfont\bfseries #4\par}
  \end{minipage}\hfill
  \begin{minipage}[c]{2.6cm}\raggedleft\includegraphics[width=2.6cm]{\fm@logo}\end{minipage}\par
  \vspace{2pt}}

% ---------- barres de secció del color del curs ----------
\renewcommand{\seccio}[2]{\par\needspace{\fmneed\baselineskip}\bigskip
  \begin{tcolorbox}[colback=curs,colframe=curs,coltext=curstext,arc=2mm,boxrule=0pt,left=3mm,right=3mm,top=1.2mm,bottom=1.2mm]
  {\large\bfseries #1\quad #2}
  \end{tcolorbox}\vspace{-2pt}}

% ---------- estrelles de dificultat del color del curs (a Batxillerat no se'n posen) ----------
\@ifundefined{estrelles}{}{%
  \renewcommand{\estrelles}[1]{{\color{curs}\ifcase#1\or\faStar\faStar[regular]\faStar[regular]\or\faStar\faStar\faStar[regular]\or\faStar\faStar\faStar\fi}}}

% ---------- capçalera i peu de les pàgines ----------
% primera pàgina: logos institucionals al peu
\fancypagestyle{primera}{\fancyhf{}\renewcommand{\headrulewidth}{0pt}
  \fancyfoot[L]{\raisebox{-0.32\height}{\includegraphics[height=0.78cm]{logo-gva.png}}\hspace{5mm}%
    \raisebox{-0.32\height}{\includegraphics[height=0.78cm]{logo-vermellar.png}}}
  \fancyfoot[R]{\scriptsize\color{muted}\href{https://creativecommons.org/licenses/by-nc-sa/4.0/deed.ca}{\ccbyncsa}\ \ CC BY-NC-SA 4.0 · fentmates.org\qquad\small\thepage}}
% la resta: el símbol a l'esquerra i el curs a la dreta, amb el títol al centre com abans
\fancyhead[L]{\raisebox{-0.3\height}{\includegraphics[height=0.62cm]{\fm@marca}}}
\fancyhead[R]{\small\bfseries\color{cursfosc}\fm@cursnom}
\renewcommand{\headrule}{\color{curs!45}\hrule height \headrulewidth}
\makeatother

%   \fmcurs{1eso}
% i, al començament del document:
%   \capcalera{Matemàtiques}{1r ESO}{Fitxa d'exercicis · Unitat 2}{Divisibilitat}
% Opcional abans de \input: \def\fmneed{14} (espai mínim abans de cada secció).
% =====================================================================
\makeatletter
\def\fm@logodir{../../logo/pdf}
\providecommand{\fmneed}{10}
\def\fm@cursnom{}

% clau del curs, color i color del text damunt del color
\newcommand{\fm@curs}[3]{\@namedef{fm@color@#1}{#2}\@namedef{fm@text@#1}{#3}}
\fm@curs{1eso}{E6B84A}{navy}
\fm@curs{2eso}{3B82F6}{white}
\fm@curs{3eso}{FF7A50}{white}
\fm@curs{4eso-a}{A78BFA}{white}
\fm@curs{4eso-b}{6D28D9}{white}
\fm@curs{1batx-mi}{9B1C1C}{white}
\fm@curs{1batx-ccss}{EF6B6B}{white}
\fm@curs{2batx-mii}{1E6B47}{white}
\fm@curs{2batx-ccss}{2BB3A3}{white}

\newcommand{\fmcurs}[1]{%
  \@ifundefined{fm@color@#1}{\PackageError{fentmates}{Curs desconegut: #1}{}}{}%
  \edef\fm@tmp{\noexpand\definecolor{curs}{HTML}{\@nameuse{fm@color@#1}}%
               \noexpand\colorlet{curstext}{\@nameuse{fm@text@#1}}}\fm@tmp
  \colorlet{cursfosc}{curs!72!black}%
  \def\fm@logo{\fm@logodir/vertical-#1.pdf}%
  \def\fm@marca{\fm@logodir/marca-#1.pdf}}

% ---------- capçalera de la primera pàgina ----------
% #1 matèria (petit), #2 curs, #3 tipus de document i unitat, #4 títol
\newcommand{\capcalera}[4]{%
  \gdef\fm@cursnom{#2}%
  \par\noindent
  \begin{tcolorbox}[colback=curs,colframe=curs,boxrule=0pt,arc=2.2mm,width=3.5cm,height=2.15cm,
      valign=center,left=3mm,right=2mm,top=1mm,bottom=1mm,nobeforeafter,box align=center]
    \color{curstext}{\footnotesize #1}\par\vspace{1pt}{\fontsize{19}{21}\selectfont\bfseries #2\par}%
  \end{tcolorbox}\hspace{5mm}%
  \begin{minipage}[c]{\dimexpr\linewidth-3.5cm-5mm-2.9cm\relax}
    {\color{cursfosc}\small\bfseries\MakeUppercase{#3}}\par\vspace{2pt}
    {\color{navy}\fontsize{20}{23}\selectfont\bfseries #4\par}
  \end{minipage}\hfill
  \begin{minipage}[c]{2.6cm}\raggedleft\includegraphics[width=2.6cm]{\fm@logo}\end{minipage}\par
  \vspace{2pt}}

% ---------- barres de secció del color del curs ----------
\renewcommand{\seccio}[2]{\par\needspace{\fmneed\baselineskip}\bigskip
  \begin{tcolorbox}[colback=curs,colframe=curs,coltext=curstext,arc=2mm,boxrule=0pt,left=3mm,right=3mm,top=1.2mm,bottom=1.2mm]
  {\large\bfseries #1\quad #2}
  \end{tcolorbox}\vspace{-2pt}}

% ---------- estrelles de dificultat del color del curs (a Batxillerat no se'n posen) ----------
\@ifundefined{estrelles}{}{%
  \renewcommand{\estrelles}[1]{{\color{curs}\ifcase#1\or\faStar\faStar[regular]\faStar[regular]\or\faStar\faStar\faStar[regular]\or\faStar\faStar\faStar\fi}}}

% ---------- capçalera i peu de les pàgines ----------
% primera pàgina: logos institucionals al peu
\fancypagestyle{primera}{\fancyhf{}\renewcommand{\headrulewidth}{0pt}
  \fancyfoot[L]{\raisebox{-0.32\height}{\includegraphics[height=0.78cm]{logo-gva.png}}\hspace{5mm}%
    \raisebox{-0.32\height}{\includegraphics[height=0.78cm]{logo-vermellar.png}}}
  \fancyfoot[R]{\scriptsize\color{muted}\href{https://creativecommons.org/licenses/by-nc-sa/4.0/deed.ca}{\ccbyncsa}\ \ CC BY-NC-SA 4.0 · fentmates.org\qquad\small\thepage}}
% la resta: el símbol a l'esquerra i el curs a la dreta, amb el títol al centre com abans
\fancyhead[L]{\raisebox{-0.3\height}{\includegraphics[height=0.62cm]{\fm@marca}}}
\fancyhead[R]{\small\bfseries\color{cursfosc}\fm@cursnom}
\renewcommand{\headrule}{\color{curs!45}\hrule height \headrulewidth}
\makeatother

\fmcurs{1batx-ccss}

\begin{document}
\thispagestyle{primera}

\capcalera{Matemáticas CCSS}{1\textsuperscript{o} Bach.}{Ficha de ejercicios · Unidad 9}{\raggedright Combinatoria y probabilidad}
\vspace{10pt}
Nombre y apellidos: \rule{6.9cm}{0.4pt}\hfill Grupo: \rule{1.6cm}{0.4pt}\hfill Fecha: \rule{2.2cm}{0.4pt}
\vspace{8pt}

\begin{tcolorbox}[colback=soft,colframe=soft,arc=2mm,left=3mm,right=3mm,top=1.5mm,bottom=1.5mm]
\small Resuelve los ejercicios en el cuaderno mostrando todos los pasos: di qué tipo de recuento usas (y por qué), dibuja el diagrama de árbol o la tabla de contingencia cuando te ayude y escribe las probabilidades como fracciones o con tres cifras decimales. Puedes comprobar los resultados con las herramientas de los apuntes interactivos.
\end{tcolorbox}

% =====================================================
\seccio{1}{Técnicas de recuento. Principio de multiplicación y factorial}

\exf{1} Un restaurante ofrece 4 primeros platos, 5 segundos y 3 postres. ¿Cuántos menús distintos (un plato de cada) se pueden elegir? ¿Cuántos, si una persona no quiere pescado y hay 2 segundos de pescado?

\ex{1} Una contraseña está formada por 3 letras (de un alfabeto de 26) seguidas de 3 cifras. ¿Cuántas contraseñas hay? ¿Cuántas no repiten ninguna letra ni ninguna cifra?

\ex{2} ¿Cuántos números de cuatro cifras (del 1000 al 9999) hay que\begin{enumerate}[label=\alph*),leftmargin=*,itemsep=2pt,topsep=2pt]\item sean pares?\item tengan las cuatro cifras distintas?\item sean capicúas (como 3443)?\end{enumerate}\recorda{empieza por la cifra que tiene condiciones; la primera cifra no puede ser 0.}

\ex{1} Calcula y simplifica:\begin{tasks}(4)
  \task $6!$
  \task $\dfrac{9!}{7!}$
  \task $\dfrac{10!}{7!\,3!}$
  \task $\dfrac{(n+1)!}{(n-1)!}$
\end{tasks}Resuelve la ecuación $\dfrac{n!}{(n-2)!}=30$.

\ex{2} ¿Cuántos números (del 1 al 999) se pueden escribir usando solo las cifras 1, 2 y 3 (que se pueden repetir)?\recorda{si los casos se separan en grupos sin elementos comunes, se suman los recuentos de cada grupo (principio de la suma).}

\ex{2} Lanzamos cuatro dados (de colores distintos). ¿Cuántos resultados hay? ¿En cuántos aparece al menos un 6?


% =====================================================
\seccio{2}{Variaciones y permutaciones}

\exf{1} Calcula:\begin{tasks}(5)
  \task $V_{8,\,3}$
  \task $VR_{5,\,3}$
  \task $P_5$
  \task $PR_6^{\,2,\,2,\,2}$
  \task $PC_5$
\end{tasks}

\ex{1} Un club de 20 socios tiene que elegir presidente, secretario y tesorero (tres personas distintas). ¿De cuántas maneras lo puede hacer?

\ex{2} Con las cifras del 1 al 9, ¿cuántos números de cinco cifras distintas se pueden formar? ¿Cuántos empiezan por 1 y acaban en 9? ¿Cuántos son pares?

\ex{2} Lanzamos una moneda 10 veces y anotamos la secuencia de caras y cruces. ¿Cuántas secuencias hay? ¿En cuántas sale exactamente 4 veces cara?

\ex{2} ¿Cuántas palabras (con sentido o sin él) se pueden formar con todas las letras de PATATA? ¿Cuántas empiezan por P? ¿Cuántas empiezan y acaban por A?

\ex{2} Queremos colocar en una estantería 4 libros de matemáticas, 3 de historia y 2 de inglés, todos distintos. ¿De cuántas maneras se puede hacer? ¿Y si los de cada materia tienen que estar juntos?

\ex{3} Siete amigos se sientan alrededor de una mesa redonda (solo importa quién tiene a cada lado). ¿De cuántas maneras lo pueden hacer? ¿Y si dos de ellos, enfadados, no quieren sentarse juntos?

\ex{2} Tres chicos y tres chicas se sientan en una fila de 6 sillas. ¿De cuántas maneras pueden hacerlo si se tienen que alternar chicos y chicas?


% =====================================================
\seccio{3}{Combinaciones y números combinatorios}

\exf{1} Calcula (usa la propiedad de los números complementarios cuando convenga):\begin{tasks}(4)
  \task $\dbinom94$
  \task $\dbinom{12}{10}$
  \task $\dbinom{20}{18}$
  \task $\dbinom70$
\end{tasks}

\ex{2} Resuelve las ecuaciones:\begin{tasks}(3)
  \task $\dbinom n2=28$
  \task $\dbinom x3=\dbinom x5$
  \task $\dbinom{10}{x}=\dbinom{10}{x+2}$
\end{tasks}\recorda{si $\binom na=\binom nb$ con $a\neq b$, entonces $a+b=n$ (números complementarios).}

\ex{2} Usa las propiedades de los números combinatorios para calcular:\begin{tasks}(3)
  \task $\dbinom94+\dbinom95$
  \task $\dbinom60+\dbinom61+\dots+\dbinom66$
  \task $\dbinom73+\dbinom74+\dbinom85$
\end{tasks}

\ex{2} Desarrolla con el binomio de Newton:\begin{tasks}(2)
  \task $(x+2)^4$
  \task $(2x-1)^4$
\end{tasks}

\ex{2} En un grupo hay 7 mujeres y 5 hombres. Hay que formar una comisión de 4 personas. ¿De cuántas maneras se puede hacer\begin{enumerate}[label=\alph*),leftmargin=*,itemsep=2pt,topsep=2pt]\item sin ninguna condición?\item si tiene que tener 2 mujeres y 2 hombres?\item si tiene que haber al menos una mujer?\end{enumerate}

\ex{1} En un examen de 10 preguntas hay que responder 7. ¿Cuántas elecciones distintas hay? ¿Y si las tres primeras son obligatorias?

\ex{2} ¿Cuántas diagonales tiene un octógono convexo? ¿Cuántos triángulos se pueden dibujar con vértices en los vértices del octógono?

\ex{2} En la Primitiva se marcan 6 números del 1 al 49. En el Euromillones se marcan 5 números del 1 al 50 y 2 «estrellas» del 1 al 12. ¿Cuántas apuestas distintas hay en cada juego?

\ex{2} Di de qué tipo es cada recuento y calcúlalo:\begin{enumerate}[label=\alph*),leftmargin=*,itemsep=2pt,topsep=2pt]\item maneras de sentarse 5 personas en 5 sillas en fila;\item maneras de elegir 3 sabores distintos (de 10) para una copa de helado;\item códigos de 4 cifras para un candado;\item maneras de repartir 3 premios distintos entre 10 participantes (cada uno gana como mucho uno);\item palabras que se pueden formar con todas las letras de RISAS.\end{enumerate}


% =====================================================
\seccio{4}{Sucesos. Probabilidad y regla de Laplace}

\exf{1} Lanzamos tres monedas. Escribe el espacio muestral y los sucesos $A=$ «al menos dos caras» y $B=$ «la primera moneda es cruz». Escribe y calcula la probabilidad de $A\cap B$, $A\cup B$, $\overline A$ y $B\cap\overline A$.

\ex{1} Lanzamos dos dados. Calcula la probabilidad de que\begin{tasks}(2)
  \task la suma sea 8;
  \task la suma sea 10 o más;
  \task el producto sea múltiplo de 3;
  \task la diferencia entre los dos sea como mucho 1.
\end{tasks}

\ex{1} Sacamos una carta de una baraja española de 40 cartas (4 palos: oros, copas, espadas y bastos; en cada palo, del 1 al 7, sota, caballo y rey). Calcula la probabilidad de sacar\begin{tasks}(4)
  \task un oro;
  \task una figura;
  \task un oro o una figura;
  \task ni un as ni una copa.
\end{tasks}

\ex{2} Una urna tiene 6 bolas blancas y 4 negras. Sacamos 3 bolas a la vez. Usa la combinatoria para calcular la probabilidad de que\begin{enumerate}[label=\alph*),leftmargin=*,itemsep=2pt,topsep=2pt]\item las tres sean blancas;\item haya exactamente dos blancas;\item haya al menos una negra.\end{enumerate}

\ex{2} Colocamos al azar cinco libros, A, B, C, D y E, en una estantería. ¿Cuál es la probabilidad de que A y B queden juntos? ¿Y de que A quede a la izquierda de B (no necesariamente al lado)?

\ex{3} En la Primitiva (6 números de 49), ¿cuál es la probabilidad de acertar los 6 números con una apuesta? ¿Y de acertar exactamente 5?

\ex{2} Se elige al azar un código de 4 cifras (del 0000 al 9999). Probabilidad de que todas las cifras sean distintas, y de que sea capicúa.

\ex{2} En una clase hay 12 chicas y 8 chicos. Se eligen 3 personas al azar. Probabilidad de que haya al menos un chico.


% =====================================================
\seccio{5}{Propiedades de la probabilidad. Probabilidad condicionada}

\exf{1} Sean $A$ y $B$ dos sucesos con $P(A)=0{,}5$, $P(B)=0{,}4$ y $P(A\cap B)=0{,}2$. Calcula:\begin{tasks}(3)
  \task $P(A\cup B)$
  \task $P(\overline A\cap\overline B)$
  \task $P(A\cap\overline B)$
  \task $P(\overline A\cup\overline B)$
  \task $P(A\,|\,B)$
  \task $P(B\,|\,A)$
\end{tasks}\recorda{leyes de De Morgan: $\overline A\cap\overline B=\overline{A\cup B}$ y $\overline A\cup\overline B=\overline{A\cap B}$.}

\ex{2} Sabemos que $P(A)=0{,}6$, $P(B)=0{,}3$ y $P(A\cup B)=0{,}7$. Calcula $P(A\cap B)$, $P(A\,|\,B)$ y $P(B\,|\,\overline A)$.

\ex{1} $A$ y $B$ son incompatibles, con $P(A)=0{,}35$ y $P(B)=0{,}45$. Calcula $P(A\cup B)$, $P(\overline A\cap\overline B)$ y $P(A\,|\,B)$.

\ex{2} En un grupo de 200 personas hay 120 mujeres. Fuman 30 mujeres y 25 hombres. Haz la tabla de contingencia y calcula la probabilidad de que una persona elegida al azar\begin{tasks}(2)
  \task fume;
  \task sea mujer y fume;
  \task fume, sabiendo que es mujer;
  \task sea mujer, sabiendo que fuma;
  \task sea hombre o fume.
\end{tasks}

\ex{2} En un instituto, el 60 \% del alumnado aprueba Matemáticas, el 50 \% aprueba Inglés y el 35 \% aprueba las dos. Elegimos un alumno al azar. Calcula la probabilidad de que\begin{tasks}(2)
  \task apruebe al menos una de las dos;
  \task no apruebe ninguna;
  \task apruebe solo una;
  \task apruebe Inglés, si ha aprobado Matemáticas;
  \task apruebe Matemáticas, si ha suspendido Inglés.
\end{tasks}

\ex{1} En una clase de 30 alumnos, 18 juegan a fútbol, 12 a baloncesto y 6 a los dos deportes. Elegimos un alumno al azar. Probabilidad de que no juegue a ninguno de los dos. Si juega a fútbol, ¿cuál es la probabilidad de que también juegue a baloncesto? ¿Y al revés?


% =====================================================
\seccio{6}{Experiencias compuestas. Sucesos y experiencias dependientes e independientes}

\exf{2} Estudia si los sucesos $A$ y $B$ son independientes en cada caso:\begin{enumerate}[label=\alph*),leftmargin=*,itemsep=2pt,topsep=2pt]\item $P(A)=0{,}5$, \ $P(B)=0{,}4$, \ $P(A\cap B)=0{,}2$\item $P(A)=0{,}3$, \ $P(B)=0{,}6$, \ $P(A\cup B)=0{,}72$\item $P(A)=0{,}4$, \ $P(B)=0{,}5$, \ $P(A\,|\,B)=0{,}6$\end{enumerate}\recorda{$A$ y $B$ son independientes si $P(A\cap B)=P(A)\cdot P(B)$ o, lo que es lo mismo, si $P(A\,|\,B)=P(A)$.}

\ex{2} Ana aprueba el carnet de conducir a la primera con probabilidad 0,7 y Pablo, con probabilidad 0,4, de forma independiente. Calcula la probabilidad de que lo aprueben a la primera\begin{tasks}(4)
  \task los dos;
  \task al menos uno;
  \task solo uno;
  \task ninguno de los dos.
\end{tasks}

\ex{2} Una urna contiene 4 bolas rojas y 6 azules. Sacamos dos bolas. Calcula la probabilidad de que las dos sean rojas y de que sean de colores distintos (a) con reemplazamiento; (b) sin reemplazamiento. En el caso (b), ¿cuál es la probabilidad de que la segunda sea roja?

\ex{2} Un jugador encesta un tiro libre con probabilidad 0,8. Hace tres tiros (independientes). Calcula la probabilidad de que enceste\begin{tasks}(4)
  \task los tres;
  \task ninguno;
  \task al menos uno;
  \task exactamente dos.
\end{tasks}

\ex{2} De una baraja española de 40 cartas sacamos tres cartas, una tras otra, sin devolverlas. Probabilidad de que las tres sean oros. Compruébalo usando combinaciones.

\ex{3} Lanzamos un dado cuatro veces. ¿Cuál es la probabilidad de obtener al menos un 6? ¿Cuántas veces lo tenemos que lanzar, como mínimo, para que la probabilidad de obtener al menos un 6 sea mayor que 0,9?

\ex{3} Una urna tiene 3 bolas blancas y 2 negras. Sacamos una bola, la devolvemos a la urna y añadimos otra bola del mismo color. Después sacamos una segunda bola. Calcula la probabilidad de que las dos sean blancas y de que la segunda sea blanca.


% =====================================================
\seccio{7}{Teorema de la probabilidad total y teorema de Bayes}

\exf{2} Una fábrica tiene tres máquinas: A hace el 50 \% de las piezas, B el 30 \% y C el 20 \%. Son defectuosas el 3 \% de las piezas de A, el 4 \% de las de B y el 5 \% de las de C. Elegimos una pieza al azar.\begin{enumerate}[label=\alph*),leftmargin=*,itemsep=2pt,topsep=2pt]\item ¿Cuál es la probabilidad de que sea defectuosa?\item Si es defectuosa, ¿cuál es la probabilidad de que la haya hecho la máquina A?\item Si no es defectuosa, ¿cuál es la probabilidad de que la haya hecho C?\end{enumerate}

\ex{2} Tenemos dos urnas: la I tiene 3 bolas blancas y 2 negras; la II, 1 blanca y 4 negras. Lanzamos una moneda: si sale cara, sacamos una bola de la I; si sale cruz, de la II. Probabilidad de que la bola sea blanca. Si ha salido blanca, ¿cuál es la probabilidad de que sea de la urna I?

\ex{3} Una enfermedad afecta al 2 \% de la población. Una prueba diagnóstica da positivo en el 95 \% de los enfermos y también, por error, en el 5 \% de las personas sanas. Si la prueba de una persona da positivo, ¿cuál es la probabilidad de que esté enferma? Si da negativo, ¿cuál es la probabilidad de que esté sana? Comenta el resultado.

\ex{2} Del alumnado de un instituto, el 30 \% va en autobús, el 50 \% a pie y el 20 \% en coche. Llegan tarde alguna vez el 10 \% de los que van en autobús, el 5 \% de los que van a pie y el 15 \% de los que van en coche. ¿Qué porcentaje del alumnado llega tarde alguna vez? Si un alumno ha llegado tarde, ¿cuál es la probabilidad de que vaya en autobús? ¿Y si no ha llegado nunca tarde, de que vaya a pie?

\ex{2} Lanzamos un dado: si sale 1 o 2, sacamos una bola de la urna A, que tiene 2 bolas rojas y 3 verdes; en caso contrario, la sacamos de la urna B, que tiene 4 rojas y 1 verde. Probabilidad de que la bola sea roja. Si la bola es verde, ¿cuál es la probabilidad de que se haya sacado de la urna A?

\vfill
\begin{center}\color{muted}\small\faCheckCircle[regular]\ \ ¡Ánimo! Comprueba tus resultados con las herramientas interactivas de los apuntes de la unidad.\end{center}

\end{document}
